% !TeX root = ../main.tex
\begingroup
\setbeamerfont{normal text}{size=\fontsize{22}{30}\selectfont}
\setbeamerfont{BIC point}{size=\fontsize{24}{30}\selectfont}
\setbeamerfont{BIC conclusion}{size=\fontsize{24}{28}\selectfont}
\begin{frame}{04 结论与展望——结论}
  \begin{SlideCanvas}
    \RedAt{55}{102}{70}{01}
    \PointAt{136}{103}{750}{近零泊松比为长度控制提供结构路径}
    \BodyAt{136}{148}{740}{已有聚合物试验与镍钛研究支持其力学潜力，结论受材料和工况约束。}
    \RuleAt{211}
    \RedAt{55}{239}{70}{02}
    \PointAt{136}{240}{750}{联合设计提供目标响应的共同定制能力}
    \BodyAt{136}{285}{740}{课题组方法提供有限变形基础，转向支架需补充径向部署与抗压评价。}
    \RuleAt{348}
    \RedAt{55}{376}{70}{03}
    \PointAt{136}{377}{750}{方案价值应由统一约束下的对照检验}
    \BodyAt{136}{422}{740}{拟比较长度误差、力曲线及局部应变，检验可行范围与性能代价。}
    \Sources{主要依据：\RefGroup；\RefWang；\RefYasmin。第17--18页为应用构想与拟议验证。}
  \end{SlideCanvas}
  \BICConclusionBand{可编程能力的应用价值在于\BICBandEmph{协调长度控制与机械支撑}}
  \note{末页回到最初的问题。可编程泊松比对轴向长度控制具有直接的机制联系，联合设计进一步处理变形需求与承载需求之间的取舍。现有论文为方法提供证据，所提出的支架工况研究仍需开展。合理的初步结果既可以是某个目标组合有收益，也可以是它在曲率、接触或材料条件下出现明确限制。谢谢老师同学。}
\end{frame}
\endgroup
